% Generated from locale/tr/content/incompleteness/arithmetization-syntax/coding-formulas.tex; do not hand-edit.


\olfileid[tr]{inc}{art}{frm}
\olsection{\printtoken{P}{formula}ın Kodlanması}

$\fn{Term}(x)$ bağıntısını ilkel özyinelemeli olarak tanımladıktan sonra,
!!{formula}s için karşılık gelen $\fn{Frm}(x)$ bağıntısını da ilkel
özyinelemeli olarak tanımlayabiliriz.

\begin{prop}
$x$'in atomik bir !!{formula} için Gödel numarası olması durumunda geçerli olan
$\fn{Atom}(x)$ bağıntısı ilkel özyinelemelidir.
\end{prop}

\begin{proof}
$x$ sayısı, ancak ve ancak aşağıdakilerden biri geçerliyse atomik bir
!!{formula} için Gödel numarasıdır:
\begin{enumerate}
\item Her $i<n$ için $\fn{Term}((z)_i)$ olacak biçimde $n,j<x$ ve $z<x$
  vardır ve $x=$
\[
\Gn{\Obj P^n_j(} \concat \fn{flatten}(z) \concat \Gn{)}.
\]
\item $\fn{Term}(z_1)$ ve $\fn{Term}(z_2)$ olacak biçimde $z_1,z_2<x$
  vardır ve $x=$
\[
\Gn{{\eq}(} \concat z_1 \concat \Gn{,} \concat z_2 \concat{\Gn{)}}.
\]
  \tagitem{prvFalse}{$x = \Gn{\lfalse}$.}{}
  \tagitem{prvTrue}{$x = \Gn{\ltrue}$.}{}
\end{enumerate}
\end{proof}

\begin{prop}
\ollabel{prop:frm-primrec}
$x$'in !!a{formula}{}ün Gödel numarası olması durumunda geçerli olan
$\fn{Frm}(x)$ bağıntısı ilkel özyinelemelidir.
\end{prop}

\begin{proof}
Bir $s$ simge dizisi, ancak ve ancak !!a{formula} ise ve $s$'nin atomik
!!{formula}s{}den kuruluş kurallarına göre nasıl kurulduğunu kaydeden,
!!{formula}s{}den oluşan bir $s_0$, \dots, $s_{k-1}=s$ kuruluş dizisi varsa
formüldür. Her $s_i$'nin
kodu (hatta $\tuple{s_0,\dots,s_{k-1}}$ dizisinin kodunun kodu), $s$'nin
$x$ kodundan küçüktür.
\end{proof}

\begin{prob}
\olref[inc][art][frm]{prop:frm-primrec} için,
\olref[inc][art][trm]{prop:term-primrec} önermesinin ilk ispatındaki çizgi
boyunca ayrıntılı bir ispat
verin.
\end{prob}

\begin{prop}
\ollabel{prop:freeocc-primrec}
Gödel numarası $x$ olan formülün $i$'inci simgesinin, Gödel numarası $z$
olan değişkenin serbest bir geçişi olması durumunda geçerli olan
$\fn{FreeOcc}(x,z,i)$ bağıntısı ilkel özyinelemelidir.
\end{prop}

\begin{proof}
Alıştırma.
\end{proof}

\begin{prob}
\olref[inc][art][frm]{prop:freeocc-primrec} önermesini ispatlayın.
!!a{formula}{}ün içinde bulunan ve kendisi de !!a{formula} olan her alt
dizenin onun bir alt-!!{formula}{}ü olduğu olgusundan yararlanabilirsiniz.
\end{prob}

\begin{prop}
$x$'in !!a{sentence}{}nin Gödel numarası olması durumunda geçerli olan
$\fn{Sent}(x)$ özelliği ilkel özyinelemelidir.
\end{prop}

\begin{proof}
Bir !!{sentence}, !!a{formula} olup !!{variable}s{}in serbest geçişlerini içermez.
Dolayısıyla $\fn{Sent}(x)$ ancak ve ancak
\begin{multline*}
\bforall{i<\len{x}}{\bforall{z<x}{}}\\
(\bexists{j<z}{z=\Gn{\Obj v_j}} \lif \lnot\fn{FreeOcc}(x,z,i))
\end{multline*}
geçerliyse geçerlidir.
\end{proof}


